% (User input window)

\paragraph{证书 Pareto 折点 \texorpdfstring{$\Longleftrightarrow$}{<->} 多分形骨架：预算前沿与谱前沿的同构}
本段把“选 \(q\) 的工程折点”与“层集谱的换相骨架”对齐为同一凸对偶对象的两种读出。

\begin{remark}[证书折点与谱骨架换相同构：同一对偶下的两类 kink]\label{rem:pom-certificate-pareto-kink-multifractal-skeleton}
对固定 \((m,\varepsilon)\)，推论 \ref{cor:pom-prime-register-q-quota-family} 给出证书族
\[
\gamma^{(q)}_{m}(\varepsilon)
:=\frac{1}{q-1}\left(\frac{1}{m}\log S_q(m)-\log 2+\frac{1}{m}\log(1/\varepsilon)\right),
\qquad q\in\ZZ_{\ge 2},
\]
其下包络
\(
\gamma^{\mathrm{cert,fin}}_m(\varepsilon)=\inf_{q\ge 2}\gamma^{(q)}_m(\varepsilon)
\)
为一族 \(1/m\) 型双曲线的分段下包络（注 \ref{rem:pom-qstar-phase-diagram-breakpoints}）。
在断点处出现两点并列最优（最优 \(q^\star\) 离散跳变），这在预算域 \(\mathbf{Bud}\) 的语言里等价于：由离散下凸闭包 \(\mathbb{T}\) 生成的正规形出现 kink（定义 \ref{def:pom-kink-set}）。
\par
另一方面，定理 \ref{thm:pom-ldp-uniform-and-wm-exact-spectrum} 与命题 \ref{prop:pom-phase-transition-kink-vs-linear-face} 表明：压力 \(\tau\) 与层集谱 \(f\) 互为 Legendre 对偶，
并且几何相变可被内部化为“最优层集不唯一”（等价于 \(\tau\) 的折点，等价于 \(f\) 的线性面）。
因此，在本文同一对偶坐标系中：
\begin{itemize}
  \item \textbf{工程侧 kink}：最优 \(q^\star\) 的跳变点，是证书下包络中“活跃仿射证书”发生换支之处；
  \item \textbf{几何侧 kink}：主导层集/相集合的换相点，是 \(\tau(q)=\sup_\gamma(f(\gamma)+q\gamma)\) 中“活跃层集”发生换支（或并列）之处。
\end{itemize}
两者由同一离散凸对偶机制强制同步：证书预算的折点集合与多分形谱的骨架换相集合在制度上同构，
从而“可控时间内的最优证书/扫描策略”与“分形边界的相变结构”并非并列现象，而是同一凸几何对象在两种接口（预算 vs.\ 谱）下的等价像。
\end{remark}

